\silentchapter{Proposisjonar}{proposisjonar}


\prop{1}{}{Formverda er alt som er tilfelle.}
\prop{1.1}{o}{Formverda er totaliteten av realiserte konfigurasjonar, ikkje av ting\footnote{Dette gjeld formelt uavhengig av om konfigurasjonen er fysisk materie, informasjonsarkitektur i ei programvare, eller eit sekvensielt tenesteforløp.}.}
\prop{1.11}{t}{Alt som har form, har nettopp denne forma og ikkje ei anna\footnote{Identiteten til eit system er uttømmande definert av konfigurasjonen til komponentane sine i det augneblikket.}.}
\prop{1.12}{t}{At ein konfigurasjon tek éin eksakt form, når uendeleg mange andre var logisk moglege, krev ei årsaksforklaring utover konfigurasjonen sjølv.}
\prop{1.121}{t}{Forklaringa ligg i relasjonen mellom den manifesterte konfigurasjonen og totaliteten av dei formene som var moglege men ikkje vart realiserte\footnote{Relasjonen konstituerer ei interaksjonsflate: det realiserte står i konstant spenning til det latente.}.}
\prop{1.1211}{t}{Dette er ein relasjonell, ikkje ein absolutt, eigenskap\footnote{Eigenskapen oppstår i differansen mellom tilstandane, analogt med korleis eit digitalt grensesnitt får form gjennom tilstandar det utelukkar.}.}
\prop{1.2}{d}{Eit \textbf{objekt\footnote{Objektet er ikkje nødvendigvis romleg; i eit distribuert nettverk er objektet ein udeleleg datapakke eller ei atomær hending.}} er den enklaste bestanddelen i ein form.}
\prop{1.21}{t}{Objektet er udeleleg og uforanderleg\footnote{Udelelegheit tyder her at vidare dekomponering øydelegg objektets kausale funksjon i den gjevne lyskjegla.}.}
\prop{1.211}{t}{Objekta utgjer formverdas substans\footnote{Substansen gjev repertoaret for kva tenester eller interaksjonar som overhovudet kan oppstå.}.}
\prop{1.2111}{t}{Sidan objektas natur ber i seg moglegheita for å inngå i konfigurasjonar, er alle moglege former allereie gjeve i og med objekta\footnote{Ein protokoll innheld potensialet for alle framtidige transaksjonar, allereie før første transaksjon er eksekvert.}.}
\prop{1.21111}{o}{Moglegheita er logisk tilgjengeleg fordi objekta er der, ikkje fordi nokon tenkjer ho\footnote{Grammatikken for grensesnittet er uavhengig av brukarintensjonen.}.}
\prop{1.22}{d}{Ein \textbf{konfigurasjon\footnote{Konfigurasjonen er reglane for samanføyning, isomorf med syntaksen i eit kodespråk.}} er ein bestemt samanheng av objekt.}
\prop{1.221}{t}{Forma er konfigurasjonens struktur\footnote{Strukturen ber informasjonen. Geometrien er berre eitt mogleg substrat for struktur.}.}
\prop{1.3}{d}{\textbf{Formrommet\footnote{Formrommet generaliserer mogleikane: det kan skildre alt frå topologien i eit API til faserommet i eit maskinlæringssamfunn.}} (\textit{morphospace}) til ein klasse er mengda av alle moglege konfigurasjonar for objekt i denne klassen.\footnote{Raup (1966); Mitteroecker \& Huttegger (2009)}}
\prop{1.31}{d}{Kvart realisert objekt er eitt eksakt punkt i dette n-dimensjonale rommet\footnote{Tilstandsrommet i eit digitalt system tilordnar kvar unik konfigurasjon ein eksakt og reversibel koordinat.}.}
\prop{1.32}{d}{Det empiriske formrommet er alltid ein projeksjon\footnote{Projeksjonen er uunngåeleg når n-dimensjonale system (som brukaropplevingar) vert reduserte til lågdimensjonale beregningar (som konverteringsratar).}.}
\prop{1.321}{t}{Inga endeleg mengd parametrar fangar den latente kompleksiteten fullt ut.\footnote{Thompson (1917)}}
\prop{1.3211}{o}{Kvar projeksjon er eit val: ho avslører nokre relasjonar og gøymer andre\footnote{Optimalisering av éin metrikk i eit system skjuler alltid degenerering av andre, unytta relasjonar.}.}
\prop{1.32111}{o}{Det finst ingen nøytral projeksjon.}
\prop{1.32112}{o}{Projeksjonen er naudsynleg for empirisk testing, men ho konstituerer ikkje forma\footnote{Grensesnittet sitt utsjånad er ein projeksjon; den underliggande algoritmens tilstand utgjer sjølve forma.}.}
\prop{1.321121}{o}{Forma eksisterer uavhengig av korleis vi måler ho.}
\prop{1.33}{d}{Klassen er den funksjonelle kategorien som definerer formrommet\footnote{Klassen definerer handlingsrommet. I tenestedesign svarer dette til domeneavgrensinga systemet opererer innafor.}.}
\prop{1.331}{t}{Ho set grensene; seleksjonstrykka formar fordelinga innanfor grensene\footnote{Brukarmønster og tekniske flaskehalsar utgjer seleksjonstrykka som modellerer frekvensen innanfor tenesta.}.}
\prop{1.4}{d}{Formrommet deler seg topologisk i tre regionar: dei busette, dei opne og dei forbodne\footnote{Topologien styrer all navigasjon; å designe eit system er å teikne grensene for desse regionane.}.}
\prop{1.41}{o}{Dei busette regionane utgjer det historiske arkivet\footnote{Dei busette regionane er dei mønstera brukarar eller agentar allereie har validert som levedyktige.}.}
\prop{1.411}{o}{Dei fungerer som ankerpunkt for all framtidig navigasjon\footnote{Eit etablert interaksjonsmønster (som eit sveip eller eit klikk) vert eit ankerpunkt for all framtidig systemutvikling.}.}
\prop{1.4111}{t}{Tyngda til eit ankerpunkt er proporsjonal med lengda på okkupasjonen og talet på uavhengige tradisjonar som har konvergert mot same posisjon\footnote{Standardiseringa av ein protokoll vert tyngre jo fleire uavhengige nodar som baserer drifta si på han.}.}
\prop{1.41111}{t}{Konvergens frå uavhengige kjelder er sterkare evidens for at ein posisjon svarer til ein reell haug i landskapet enn konvergens frå éin tradisjon åleine.}
\prop{1.42}{o}{Dei opne regionane utgjer det \textit{tilstøytande moglege\footnote{Moglegheitsrommet i oppdateringar av programvare ligg utelukkande i desse tilstøytande regionane.}}.}
\prop{1.421}{o}{Dei er tilgjengelege men uaktualiserte\footnote{Latente funksjonar i eit system er opne, men krev ein utløysande agent for aktualisering.}.}
\prop{1.4211}{t}{Det tilstøytande moglege er ikkje uendeleg.}
\prop{1.42111}{t}{Det er avgrensa av naboskapet til dei allereie busette regionane og av grensene til dei forbodne regionane\footnote{Arkitekturens tekniske gjeld og domenekrava dikterer rigorøst kva for greiner av kodebasen som kan utforskast vidare.}.}
\prop{1.42112}{o}{Naboskapet til dei busette regionane trekkjer til seg; grensene til dei forbodne støyter frå\footnote{Attraktoren er brukarverdien; den frastøytande krafta er kognitiv friksjon eller teknisk latens.}.}
\prop{1.421121}{o}{Det realiserte rommet er ein funksjon av begge.}
\prop{1.43}{d}{Dei forbodne regionane er dei som vert utelukka av materielle, strukturelle eller energetiske avgrensingar under dei gjeldande vilkåra\footnote{I informasjonssystem er dei forbodne regionane ofte definerte av minnekapasitet, tryggleikspolitiske reglar eller tidsavbrot.}.}
\prop{1.431}{t}{Det som er forbode under eitt sett av vilkår, treng ikkje vere forbode under eit anna\footnote{Ein overgang til asynkron handtering gjer umiddelbart forbodne arkitekturar tilgjengelege og opne.}.}
\prop{1.4311}{o}{Grensa mellom det opne og det forbodne er dynamisk.}
\prop{1.43111}{o}{Eit materiale som gjer ein region forboden, kan verte erstatta av eit materiale som opnar same region.}
\prop{1.43112}{o}{Grensa er ikkje ein eigenskap ved forma; ho er ein eigenskap ved vilkåra.}

